% ============================================================================
%  KCGL ERP - Development Documentation
%  Module: PMS (Performance Management System)
%  Template: same layout as the MOM Tracking document (MOM.pdf)
%  Compile with pdfLaTeX on Overleaf. No external assets required - screenshots
%  are optional and drop into ./screens/ (see \screenshot below).
% ============================================================================
\documentclass[11pt,a4paper]{article}

\usepackage[utf8]{inputenc}
\usepackage[T1]{fontenc}
\usepackage[a4paper,margin=1in,headheight=15pt,headsep=18pt]{geometry}
\usepackage{fancyhdr}
\usepackage{graphicx}
\usepackage{float}
\usepackage{booktabs}
\usepackage{tabularx}
\usepackage{longtable}
\usepackage{array}
\usepackage{enumitem}
\usepackage{xcolor}
\usepackage{amssymb}
\usepackage{caption}
\usepackage[hidelinks]{hyperref}

% ---- page furniture --------------------------------------------------------
\definecolor{kcglblue}{HTML}{2F3192}

\pagestyle{fancy}
\fancyhf{}
\fancyhead[L]{\small Kala Care ERP}
\fancyhead[R]{\small PMS -- Performance Management System}
\fancyfoot[C]{\thepage}
\renewcommand{\headrulewidth}{0.4pt}

\fancypagestyle{plain}{%
  \fancyhf{}\fancyfoot[C]{\thepage}\renewcommand{\headrulewidth}{0pt}}

\setlength{\parindent}{0pt}
\setlength{\parskip}{6pt}

\captionsetup{font=small,labelfont=bf}

% ---- helpers ---------------------------------------------------------------
% Requirement list: R1., R2., ... with a hanging indent (same as MOM.pdf)
\newlist{reqlist}{enumerate}{1}
\setlist[reqlist]{label=\textbf{R\arabic*.},leftmargin=3.1em,itemsep=6pt,
                  topsep=6pt,labelsep=0.6em}

% Bulleted lists
\setlist[itemize]{leftmargin=1.6em,itemsep=3pt,topsep=4pt}

% A screenshot that degrades to a framed placeholder when the image is absent,
% so the document always compiles. Drop the PNG into ./screens/ to fill it in.
\newcommand{\screenshot}[2]{%
  \begin{figure}[H]\centering
  \IfFileExists{#1}%
    {\includegraphics[width=\linewidth,keepaspectratio]{#1}}%
    {\setlength{\fboxsep}{0pt}%
     \fbox{\parbox[c][42mm][c]{0.98\linewidth}{\centering\small\itshape
       Screenshot placeholder\\[2pt]\normalfont\ttfamily #1}}}%
  \caption{#2}\end{figure}}

% Section-opening helpers, worded exactly like the MOM document
\newcommand{\covers}[1]{\textit{Covers requirements #1.}\par}
\newcommand{\purpose}[1]{\textbf{Purpose.} #1\par}
\newcommand{\devproc}{\textbf{Development process.}\par}
\newcommand{\handle}[1]{\textbf{How to handle.} #1\par}
\newcommand{\builtfrom}[1]{\textbf{Built from.} #1\par}

% Quoted text (avoids raw backtick quotes in the source)

\newcommand{\qt}[1]{\textquotedblleft #1\textquotedblright}

% Shorthands used throughout the tables
\newcommand{\file}[1]{\texttt{\small #1}}
\newcommand{\col}[1]{\textsc{\small #1}}

\begin{document}

% ============================== TITLE =======================================
\begin{center}
  {\LARGE\bfseries PMS -- Performance Management System}\\[10pt]
  {\large Development Documentation, Kala Care ERP System}\\[14pt]
  {\normalsize Prepared by: Vishal Chavan \& Nikhil Reddy \quad$|$\quad Date: August 2026}
\end{center}

\vspace{10pt}

% ============================== 1. INTRO ====================================
\section{Introduction}

This document describes the \textbf{PMS (Performance Management System)} module of the
Kala Care ERP. PMS is the module that turns the KOEL and ERP data exports into the
management sheets the business actually reviews: branch sales against the annual
operating plan (AOP), service-engineer productivity, SR allocation and closure,
the training master, and the yearly reports -- Service Penetration, AMC \& Bandhan
Projection, Customer Delight Index and Service Load \& Response.

Everything PMS prints is either \emph{counted} from an uploaded file or \emph{asserted}
in the AOP \& Master page. The module never invents a figure: a target that nobody has
entered prints a dash, and a value that no upload can produce (last year's closed
actual, FTR / FVR) is typed by a person and marked as such. Where a file's own column
disagrees with a value the module could recompute, the file's column wins and the
reason is recorded.

The module is organised into the following pages:
\begin{itemize}
  \item \textbf{AOP \& Master} -- every target and every mapping table PMS reads (9 tabs).
  \item \textbf{Sales \& Labour Report} -- Spare and Labour sale against the monthly AOP.
  \item \textbf{Employee Productivity} -- service-engineer-wise SR closure, leads and CDI.
  \item \textbf{SR Allocation Report} -- SRs allocated vs closed, engineer and date wise.
  \item \textbf{Training Report} -- the engineers' skill / training master.
  \item \textbf{Annual Reports} -- one page holding the four yearly sheets.
\end{itemize}

The document is organised the same way the MOM Tracking document is: first the roles
and the requirements received from the client, then the data model and the source
files, then the development explained \textbf{page by page and report by report} --
each section naming the source file, the exact columns that report reads, the files in
the codebase responsible for it, and the screenshot of the delivered screen.

% ============================== 2. ROLES ====================================
\section{Roles \& Access}

PMS is not a role, it is a \emph{right}. Any user can be given \qt{PMS Access} from
Profile $\rightarrow$ Edit Employee; the Master Admin always has it. The AOP \& Master
page is narrowed further, \textbf{per tab}, so a user can be given (say) only the CDI
Target tab, and only in read-only mode.

\begin{table}[H]\centering\small
\begin{tabularx}{\linewidth}{@{}p{3.4cm}X@{}}
\toprule
\textbf{Role / Right} & \textbf{Capabilities} \\
\midrule
Master Admin &
Full access to every PMS page and every AOP \& Master tab; uploads and clears data;
manages the SE UID Master on the Profile page; grants PMS rights to other users. \\
\addlinespace
PMS Access (any role) &
Opens the PMS pages granted to them -- Sales \& Labour, Employee Productivity,
SR Allocation, Training Report, Annual Reports. AOP \& Master opens only if at least
one tab is granted. \\
\addlinespace
AOP right: \textit{view} &
The granted tabs open read-only: the save / sync / reset / master buttons are hidden
and every cell is locked. Enforced again on the server, per tab. \\
\addlinespace
AOP right: \textit{edit} &
The granted tabs can be edited and saved. \\
\addlinespace
AOP rights admins &
The user ids listed in \file{VITE\_PMS\_AOP\_ADMIN\_IDS} /
\file{AOP\_RIGHTS\_ADMIN\_IDS} always hold every tab at \textit{edit}. \\
\addlinespace
Export &
Excel export is available to the Master Admin always, and to other users when the
\file{can\_export} permission is enabled on their account. \\
\bottomrule
\end{tabularx}
\caption{Role- and right-based access for the PMS module}
\end{table}

The nine AOP \& Master tabs that rights are granted on are: \textit{Target Master},
\textit{SR Type Master (Sales and Labour)}, \textit{SR Type Master (MaxTTR)},
\textit{SR Type Master (EFSR)}, \textit{Lead Category Master}, \textit{CDI Target
Master}, \textit{AMC \& Bandhan AOP}, \textit{SR Type Master (Service Load)} and
\textit{Service Load AOP}. The same nine keys exist in three places and must stay in
step: \file{client/src/utils/pagePermission.js} (\texttt{AOP\_TABS}),
\file{server/app/controllers/user\_controller.py} (\texttt{AOP\_TAB\_KEYS}) and the
\texttt{tab=} argument every AOP endpoint passes to \texttt{\_require\_aop} in
\file{server/app/routes/pms\_routes.py}.

% ============================== 3. REQUIREMENTS =============================
\section{Client Requirements}

The following requirements were received from the client and finalised for the PMS
module. They are referenced by number in the page-wise sections that follow.

\begin{reqlist}

\item \textbf{Accumulating Spare / Labour upload.} Two Excel files feed the sales side --
the \emph{Part Sale} (spares) file and the \emph{Labour Revenue} file. Rows accumulate
across daily / weekly / monthly extracts, and a re-upload of the same line must never
double-count: identity is \col{claim invoice number} plus the file's line discriminator
(\col{part number} for spares, \col{sr number} + \col{labour part num} for labour). A
changed line updates in place, an identical line is reported as a duplicate, and a row
with no invoice number falls back to a full-content hash.

\item \textbf{Flexible column matching.} Headers are matched on their alphanumeric
skeleton, so \qt{Claim Invoice Date} and \qt{CLAIM  INVOICE DATE.} are one column. Only
three columns are critical -- \col{branch id}, \col{claim invoice date} and
\col{net taxable amount}; a file missing any of them is rejected with the missing names,
and every unrecognised column is preserved rather than dropped.

\item \textbf{Cancelled invoices.} An invoice can be cancelled from the preview. The row
stays in the database as an audit trail but is excluded from every generated report, and
can be restored.

\item \textbf{AOP Target Master.} Branch-wise \emph{Spare} and \emph{Labour} targets for a
whole financial year (April--March) on one grid, entered and displayed in Lakh and stored
in rupees. The responsible person shown on the report is the branch's Branch Admin, and
the branch list follows the ERP's own branch master.

\item \textbf{Working-days calendar.} A month's target is spread over its \emph{working}
days, never its calendar days: Sundays are off, and the actual holidays are ticked on a
calendar. Maharashtra and Karnataka keep separate calendars, so a region's branches are
measured on their own days. A month with nothing ticked falls back to a typed count, and
that in turn to a universal per-calendar-month master.

\item \textbf{Sales and Labour report for any period.} Month-to-date by default, and any
week, month, quarter or year on the period picker. For every branch the report shows the
period Target, the Daily Target, the sale on the last day, the prorated Target till date,
the Achievement till date, the invoice count, \% Achieved, Short-Fall and the Balance for
the month.

\item \textbf{Independent as-on dates.} The Part Sale and the Labour files usually end on
different dates. Each side must be measured against the targets up to \emph{its own} last
data date, so a labour file that is a week behind does not read as a shortfall.

\item \textbf{Breakdown tables.} The same period broken down Region-wise (with its AOP
target), Segment-wise, Service-Report-Type-wise (through the SR Type Master) and
Category-wise (spares only, with quantity and line items).

\item \textbf{Branch-wise report.} One or more branches picked by hand, with Week-wise,
Month-wise, Segment-wise, Service-Report-Type-wise and Category-wise detail computed from
only those branches' rows.

\item \textbf{Whole-year report.} FY, Quarterly and Month-wise target-vs-achievement for
the full financial year, independent of the period picker, including the run-rate needed
to close the year.

\item \textbf{SR Type to Head masters.} The raw Service Report Types of the uploaded files
are grouped into reporting heads. Four separate masters are needed because four files
carry different vocabularies and two reports group the \emph{same} file differently --
Sales and Labour, MaxTTR, EFSR and Service Load. Values are synced out of the uploaded
data, mapped by hand, and an unmapped type must remain visible as an \qt{Unmapped} bucket
rather than being silently merged.

\item \textbf{SE UID Master.} The files identify a service engineer in two different ways
-- by \col{se name} (MaxTTR, CDI) and by \col{service engineer uid} (LMS, EFSR). One
master, maintained by hand or by Excel import on the Profile page, bridges the two,
matching on the squashed upper-case name so spacing and case never split one engineer into
two. It also holds the branch used to pin an engineer no file can place.

\item \textbf{Employee Productivity report.} Per service engineer: SRs closed (split by SR
Type head), working days, days present, productivity, CDI feedback, leads raised (split by
product category) and the spare / labour conversion amounts -- rolled up by branch, branch
group, region and overall.

\item \textbf{Attendance on task end, not close.} Days present must count the days the
engineer finished the job in the field (\col{sr task end date}), not the day the office
closed the SR -- the latter lands late and in batches.

\item \textbf{Lead attribution through the master only.} A lead counts for an engineer
because its \col{service engineer uid} is that engineer's UID \emph{in the SE UID Master}.
A UID no master row carries must be reported as an unlinked lead, not guessed.

\item \textbf{SR Allocation report.} From the EFSR Report: SRs \emph{allocated} on
\col{task assigned date} and SRs \emph{closed} on \col{task end date}, engineer wise, with
the date columns bucketed by day, week or month, plus the branch-wise open SR load.

\item \textbf{Training Report.} The KOEL Training Report export imported as a master --
eight fixed columns and every other column dynamic -- searchable two ways: by employee
(their skills and full record) and by skill (everyone who holds it).

\item \textbf{Annual: Service Penetration.} How much of the installed base was touched by
a closed service call -- IND and PG population from the asset master, unique assets with a
closed SR from the MaxTTR file, and the penetration \%.

\item \textbf{Annual: AMC and Bandhan Projection.} Per branch: last year's actual, this
year's AOP projection and the best month (all typed), the live AMC population and the
month's payments (both counted).

\item \textbf{Annual: AMC monthly sheet.} The same AMC file read by agreement
\emph{category} instead of by branch -- Bandhan (with its MH / KAR split), KALA AMC, KOEL
Corporate AMC (cumulative), expiries, renewals and Live AMC -- one financial year wide,
with the running month opened up into its weeks.

\item \textbf{Annual: Customer Delight Index.} Promotor / Passive / Detractor feedback per
branch, scored as (Promotor $-$ Detractor) $\div$ all $\times$ 100, against the CDI AOP
target, with the two previous financial years as cumulative columns.

\item \textbf{Annual: Service Load and Response.} SR load by type and by branch,
productivity (calls per person per day), the 4-hour response and the 24 / 48-hour closure
compliance, and the FTR / FVR figures -- as six tabs of one sheet.

\item \textbf{Typed figures where nothing can be counted.} FTR \%, FVR \%, last year's
closed actual and every AOP target are entered by a person in AOP and Master and printed
exactly as entered. An unset period shows a dash, which is not the same statement as 0\%.

\item \textbf{Excel export everywhere.} Every report exports to a formatted workbook that
mirrors the screen -- the ERP blue band, the header row, region totals and the overall row
-- respecting the \file{can\_export} permission. Dates export as real Excel dates so date
filters keep working.

\item \textbf{Per-user PMS rights.} PMS access is granted per user, and the AOP and Master
page per tab (none / view / edit), enforced on both the screen and the server.

\item \textbf{Performance over a remote database.} The database sits across the network,
so every report aggregates \emph{in SQL} and returns a few thousand grouped records at
most; results are cached against a data fingerprint and re-used until an upload changes
it. The payloads stay raw per-day, so changing the period, the granularity or a filter
re-aggregates in the browser without another round trip.

\item \textbf{Branch identity discipline.} Branches are resolved by \col{branch id}, never
by name -- files spell the same branch differently and the spelling changes between
imports. A branch code no KALA master knows (another dealer's code, which the EFSR file
carries on about 1\% of its rows) must fold into a single visible \qt{Unmapped Branch}
bucket rather than growing a ghost branch row.

\item \textbf{Consistent India-local timestamps.} Every stored timestamp is India local
time, so a date reads the same for every user whatever the server or browser timezone.

\end{reqlist}

% ============================== 4. SOURCE FILES =============================
\section{Source Files -- What Feeds PMS}

PMS reads two kinds of input. The \emph{Sales} files are uploaded on the Sales and Labour
page itself and the Training file on the Training Report page; every other file is
uploaded on the ERP's \textbf{Import} page and is shared with the rest of the ERP.

\begin{table}[H]\centering\small
\begin{tabularx}{\linewidth}{@{}p{4.5cm}p{3.3cm}X@{}}
\toprule
\textbf{File} & \textbf{Uploaded on} & \textbf{Which PMS reports read it} \\
\midrule
Part Sale (spares) & Sales and Labour & Sales and Labour report, all views \\
Labour Revenue & Sales and Labour & Sales and Labour report, all views \\
Training Report & Training Report & Training Report (By Employee / By Skill) \\
\addlinespace
Response Time \& MaxTTR Details & Import page & Employee Productivity; SR Allocation (home
branch); Service Penetration (serviced assets); Service Load and Response \\
EFSR Report & Import page & SR Allocation; Employee Productivity (Allocate SR) \\
LMS Data for ERP & Import page & Employee Productivity (leads and conversion amounts) \\
CDI Detail Report & Import page & Employee Productivity (CDI); Annual CDI sheet \\
Asset Detailed Report & Import page & Service Penetration (population) \\
AMC Population Report & Import page & AMC and Bandhan Projection; AMC monthly sheet \\
Anubandhan Plus / Anubandhan / BandhanPlus / Regular Bandhan quotes &
Import page & AMC and Bandhan Projection (the month's payments) \\
Open SR Load Report & Import page & SR Allocation (branch-wise open SR) \\
\bottomrule
\end{tabularx}
\caption{The uploaded files behind the PMS reports}
\end{table}

Two rules apply to every file. Dates are read \emph{date-only} and never parsed
day-first when the value is year-first, so \file{2026-07-01} can never become 7 January.
Branch ids are normalised to the underscore form used across the ERP, so
\file{420435-1}, \file{420435 1} and \file{420435\_1} are one branch.

% ============================== 5. DATA MODEL ===============================
\section{Data Model (Overview)}

Every table below lives in \file{server/app/models/pms\_model.py}. The rule that shapes
the whole model: \emph{counts} are stored per branch and rolled up, while
\emph{percentages and ratios} need a target of their own at branch, region and company
level, because a percentage cannot be summed.

\begin{longtable}{@{}p{5.2cm}p{10.0cm}@{}}
\caption{Key entities of the PMS module}\\
\toprule
\textbf{Table} & \textbf{Purpose} \\
\midrule
\endfirsthead
\multicolumn{2}{@{}l}{\small\itshape Table \thetable{} -- continued}\\
\toprule
\textbf{Table} & \textbf{Purpose} \\
\midrule
\endhead
\bottomrule
\endfoot
\file{pms\_sales\_records} & Every invoice line of the Part Sale and Labour files, with
the dedupe key, the cancelled flag and the full standard columns as real columns. \\
\file{pms\_upload\_batches} & One row per upload -- file name, total / inserted / updated
/ duplicate / skipped rows, who and when. Shared with the Training import. \\
\file{pms\_branch\_targets} & AOP Master Target tab: one row per (month, branch) with the
Spare and Labour target, region and responsible person. \\
\file{pms\_month\_settings} & A month's working days, per region, with a universal
\file{ALL-MM} fallback row. \\
\file{pms\_holidays} & One row per (date, region) -- the days that carry no target. \\
\file{pms\_heads}, \file{pms\_sr\_type\_map} & SR Type to Head master for the Sales and
Labour report. \\
\file{pms\_maxttr\_heads}, \file{pms\_maxttr\_sr\_type\_map} & The same for the MaxTTR
file (Employee Productivity). \\
\file{pms\_efsr\_heads}, \file{pms\_efsr\_sr\_type\_map} & The same for the EFSR file
(SR Allocation, Allocate SR). \\
\file{pms\_service\_load\_heads}, \file{pms\_service\_load\_sr\_type\_map} & The same
MaxTTR column grouped a second, different way for the Service Load sheet. \\
\file{pms\_lead\_categories}, \file{pms\_lead\_raised\_for\_map} & Lead Raised For to
product category -- the Product Wise Lead Count columns. \\
\file{pms\_se\_uid\_master} & Service engineer NAME to UID (and an optional pinned
branch); the bridge between the name-keyed and UID-keyed files. \\
\file{pms\_cdi\_targets} & The CDI AOP \%, per (FY, scope, key) -- branch, region or
overall. \\
\file{pms\_amc\_targets} & Per (FY, branch): last year's actual, the year's projection and
the best month -- all typed. \\
\file{pms\_amc\_category\_targets} & The AOP of each row of the AMC monthly sheet, whose
rows are agreement categories rather than branches. \\
\file{pms\_service\_load\_targets} & The monthly SR-closure target per branch -- the one
figure every Service Load rollup is derived from. \\
\file{pms\_service\_load\_pct\_targets} & The AOP of the Service Load sheet's percentage
and ratio rows, per (FY, metric, scope, key). \\
\file{pms\_service\_load\_manual} & The typed FTR / FVR percentages, per metric and
period (a month or a financial year). \\
\file{pms\_training\_records} & The Training Report master -- eight fixed columns plus
every other column of the file kept verbatim. \\
\end{longtable}

% ============================== 6. PAGE-WISE ================================
\section{Page-wise Development Process}

This section explains, page by page and report by report, what each screen does, which
client requirements it covers, \textbf{which file and which columns it reads}, which files
in the codebase are responsible for it, and the screenshot of the delivered screen.

% ---------------------------------------------------------------------------
\subsection{AOP and Master -- Target Master}
\covers{R4, R5, R27}

\purpose{The annual operating plan itself: every branch's Spare and Labour target for a
whole financial year, and the calendar those targets are spread over.}

\builtfrom{Screen \file{client/src/pages/AOPMaster.jsx} (tab \texttt{targets}).
API \texttt{GET/POST /pms/targets/year}, \texttt{GET/POST /pms/holidays}.
Server \file{pms\_controller.py} -- \texttt{get\_targets\_year\_payload},
\texttt{save\_targets\_year}, \texttt{list\_holidays}, \texttt{save\_holidays}.
Tables \file{pms\_branch\_targets}, \file{pms\_month\_settings}, \file{pms\_holidays}.}

\devproc
\begin{itemize}
  \item The grid is a whole financial year at once, April to March, one row per branch and
  two rows of inputs per branch (Spare, Labour), with quarter and year totals computed as
  the user types.
  \item Targets are \emph{entered and shown in Lakh} and stored in rupees, which is what
  the report arithmetic uses. A zero target comes back as an empty cell rather than a wall
  of noughts.
  \item The branch list, the branch names and the responsible person are taken from the
  ERP's own branch master and the branch admins, so the AOP grid and the rest of the ERP
  always name a branch identically. The AOP spelling is then treated as canonical by every
  PMS report.
  \item Working days are set per month \emph{and per region}: Maharashtra and Karnataka
  have different holidays. Ticking the actual dates on the calendar makes any part-period
  exact -- the working days of a range are the days in it, minus Sundays, minus the dates
  ticked for that region.
  \item Precedence is deliberate: a month with dates ticked uses the calendar (it knows
  \emph{which} days, not just how many); a month with nothing ticked uses its own typed
  count; failing that the universal \file{ALL-MM} master; failing that, all days except
  Sundays.
\end{itemize}

\handle{Pick the financial year, type the Spare and Labour targets in Lakh, open the month
calendar to tick holidays for MH / KA, then Save. Every Sales and Labour figure --
Target, Daily Target, Target till date, \% Achieved -- follows from this tab.}

\screenshot{screens/aop-target-master.png}{AOP and Master -- Target Master: the financial-year
grid and the working-days calendar}

% ---------------------------------------------------------------------------
\subsection{AOP and Master -- the four SR Type Masters}
\covers{R11}

\purpose{Group the raw Service Report Types that arrive in the uploaded files into the
reporting heads each report prints.}

\builtfrom{Screen \file{AOPMaster.jsx} (tabs \texttt{srtypes}, \texttt{mxtypes},
\texttt{eftypes}, \texttt{sltypes}).
API \texttt{/pms/sr-types}, \texttt{/pms/maxttr-sr-types}, \texttt{/pms/efsr-sr-types},
\texttt{/pms/service-load-sr-types} (each with \texttt{/sync}, \texttt{/reset} and a heads
endpoint). Server \file{pms\_controller.py}. Tables as listed in Section 5.}

There are four masters and not one, and the reason matters:

\begin{table}[H]\centering\small
\begin{tabularx}{\linewidth}{@{}p{3.6cm}p{3.5cm}X@{}}
\toprule
\textbf{Tab} & \textbf{Reads which column} & \textbf{Heads} \\
\midrule
SR Type Master\newline (Sales and Labour) &
Part Sale \col{claim invoice sr type}; Labour \col{sr type} &
Warranty, Post Warranty, AMC, KOEL AMC, OTC Order \\
\addlinespace
SR Type Master\newline (MaxTTR) &
MaxTTR \col{sr type} &
Warranty, Post Warranty, AMC, KOEL AMC, Courtesy Visit, Others \\
\addlinespace
SR Type Master\newline (EFSR) &
EFSR \col{sr type} &
the same six; the EFSR file adds Commercial and RECD Kit \\
\addlinespace
SR Type Master\newline (Service Load) &
MaxTTR \col{sr type}, grouped a \emph{second} way &
CSP, Post Warranty, Warranty, KOEL AMC, Dealer AMC, Courtesy Visit, Others \\
\bottomrule
\end{tabularx}
\caption{Why four SR Type masters exist}
\end{table}

\devproc
\begin{itemize}
  \item Each tab lists the distinct SR Types found in \emph{its own} file --
  \textbf{Sync} rescans the uploaded data and adds any new value, so a type that appears
  for the first time is never silently lost.
  \item \textbf{Reset} restores the defaults the business gave; the live mapping the
  reports read is always the database table, and the built-in defaults only seed it.
  \item A head can be added or deleted from the Head master beside the list, so the
  reporting buckets can change without a code change.
  \item An SR Type with no head is \emph{not} dropped: it lands in a visible
  \qt{Unmapped} bucket, so the head breakdown always adds up to the total and the gap is
  reported instead of hidden.
  \item Sales and Labour, and Service Load, group the same MaxTTR column differently on
  purpose: Employee Productivity folds CSP into Warranty and Dealer AMC into AMC, while
  the Service Load sheet prints CSP and Dealer AMC as rows of their own. Sharing one
  master would force one of the two reports to lie.
\end{itemize}

\handle{Open the tab, press Sync after a new upload, give every unmapped type its head,
Save. The change takes effect on the next report generation.}

\screenshot{screens/aop-sr-type-master.png}{AOP and Master -- SR Type Master: the synced
types and their heads}

% ---------------------------------------------------------------------------
\subsection{AOP and Master -- Lead Category Master}
\covers{R13, R15}

\purpose{Turn the LMS file's \col{lead raised for} values into the columns of the
Employee Productivity report's \emph{Product Wise Lead Count} group.}

\builtfrom{Screen \file{AOPMaster.jsx} (tab \texttt{leadcats}). API
\texttt{/pms/lead-categories} with \texttt{/sync}, \texttt{/reset}, \texttt{/lead-cats}.
Server \file{pms\_controller.py} -- \texttt{list\_lead\_map}, \texttt{save\_lead\_map},
\texttt{sync\_lead\_map}. Tables \file{pms\_lead\_categories},
\file{pms\_lead\_raised\_for\_map}.}

\devproc
\begin{itemize}
  \item The category list is the report's column list, in master order -- Allied Oil,
  Allied Product, AMC, Battery, Coolant, DEF, K-Oil, Others, Spares, Whole Goods, plus
  anything added later.
  \item Every distinct \col{lead raised for} value in the LMS upload is synced into the
  mapping table and given a category by hand (for example \emph{BD Spares} to
  \emph{Spares}, \emph{New DG} to \emph{Whole Goods}).
  \item A value with no category lands in an \qt{Unmapped} column, so the product split
  always adds up to the lead count.
\end{itemize}

\handle{Sync after an LMS upload, map the new values, Save.}

\screenshot{screens/aop-lead-category-master.png}{AOP and Master -- Lead Category Master}

% ---------------------------------------------------------------------------
\subsection{AOP and Master -- CDI Target Master}
\covers{R21, R23}

\purpose{The AOP column of the Annual Reports' Customer Delight Index sheet -- one
percentage per financial year and per row of that report.}

\builtfrom{Screen \file{AOPMaster.jsx} (tab \texttt{cditargets}). API
\texttt{GET/POST /pms/cdi-targets}. Server \texttt{list\_cdi\_targets},
\texttt{save\_cdi\_targets}. Table \file{pms\_cdi\_targets}.}

\devproc
\begin{itemize}
  \item A CDI score is a percentage, and percentages cannot be added up, so the target is
  keyed on a \emph{scope} pair rather than on a branch: \texttt{branch} with the branch
  id, \texttt{region} with MH or KA, and \texttt{overall} with \texttt{ALL} for the KCGL
  Overall row.
  \item The printed sheet sets the target at region and overall level only; branch targets
  are optional and stay empty until somebody fills them in. Nothing is derived from
  anything else -- an unset row simply shows no AOP.
\end{itemize}

\handle{Pick the financial year, type the target \% on the rows the business sets, Save.}

\screenshot{screens/aop-cdi-target.png}{AOP and Master -- CDI Target Master}

% ---------------------------------------------------------------------------
\subsection{AOP and Master -- AMC and Bandhan AOP}
\covers{R19, R20, R23}

\purpose{The figures of the two AMC sheets that the business \emph{asserts} rather than
counts. The tab carries two tables, because the two sheets have different rows.}

\builtfrom{Screen \file{AOPMaster.jsx} (tab \texttt{amctargets}). API
\texttt{GET/POST /pms/amc-targets}. Server \texttt{list\_amc\_targets},
\texttt{save\_amc\_targets}. Tables \file{pms\_amc\_targets} (per branch),
\file{pms\_amc\_category\_targets} (per agreement category).}

\devproc
\begin{itemize}
  \item \textbf{Per branch} (first table) -- three columns of the AMC and Bandhan
  Projection sheet: \emph{F26 ACT D/BAMC} (last year's actual), \emph{F27 PROJ AOP D/BAMC}
  (this year's projection) and \emph{BEST ACT AOP D/BAMC (M)} (the year's best month).
  All three come from here and nowhere else; an empty cell prints a dash rather than a
  figure derived behind the reader's back.
  \item Last year's actual is typed for a concrete reason: the AMC Population Report is a
  \emph{snapshot} keyed on instance id -- one row per genset carrying only its latest
  agreement -- so a renewal overwrites the agreement it replaced and a closed year's count
  shrinks every month. Who set that figure, and when, is stamped on the row.
  \item \textbf{Per agreement category} (second table) -- the AOP of each row of the AMC
  monthly sheet, whose rows are categories (KOEL Bandhan, KALA AMC, KOEL Corporate AMC,
  expiries, renewals, Live AMC) and therefore cannot take a sum of branch targets.
  \item Region and company rows need no target of their own here: these are counts, so the
  report adds its branches up. (The CDI target is a percentage, which is exactly why that
  one needs a target per row.)
\end{itemize}

\handle{Pick the financial year, fill the branch table and the category table, Save. The
figures appear on the Annual Reports' two AMC tabs.}

\screenshot{screens/aop-amc-bandhan.png}{AOP and Master -- AMC and Bandhan AOP: the branch
table and the category table}

% ---------------------------------------------------------------------------
\subsection{AOP and Master -- Service Load AOP}
\covers{R22, R23}

\purpose{Every target the Service Load and Response sheet is measured against, plus the
two rows nothing uploaded can produce.}

\builtfrom{Screen \file{AOPMaster.jsx} (tab \texttt{sltargets}). API
\texttt{GET/POST /pms/service-load-targets}, \texttt{GET/POST /pms/service-load-manual}.
Server \texttt{list\_service\_load\_targets}, \texttt{save\_service\_load\_targets},
\texttt{list\_service\_load\_manual}, \texttt{save\_service\_load\_manual}. Tables
\file{pms\_service\_load\_targets}, \file{pms\_service\_load\_pct\_targets},
\file{pms\_service\_load\_manual}.}

\devproc
\begin{itemize}
  \item \textbf{SR closure target} -- one number per (month, branch). Everything else on
  that sheet rolls up from it: a branch row is its twelve months summed, a region row is
  those months across the region's branches, the Total row is the same across every
  branch, and the sheet's monthly \emph{AOP Target} strip is one month across every
  branch. One table, every rollup consistent by construction.
  \item \textbf{Percentage and ratio targets} -- Productivity (Calls PP PD), 4 HRS
  RESPONSE, SR CLOSED WITHIN 24 HRS and SR CLOSED WITHIN 48 HRS. These cannot be summed,
  so each takes a value per (FY, metric, scope, key), the same scope idea the CDI target
  uses.
  \item \textbf{FTR / FVR} -- typed per month and per cumulative financial year, with
  absolute keys, so a figure entered for FY25--26 keeps meaning FY25--26 when the report
  moves on to the next year. A period nobody has filled in shows a dash, which is not the
  same statement as 0\%.
\end{itemize}

\handle{Enter the monthly SR-closure targets per branch, the percentage AOPs at region /
overall level, and the FTR / FVR figures the business supplies. Save.}

\screenshot{screens/aop-service-load.png}{AOP and Master -- Service Load AOP}

% ---------------------------------------------------------------------------
\subsection{Profile -- SE UID Master (Master Admin only)}
\covers{R12, R15, R27}

\purpose{The bridge between the files that name a service engineer and the files that
number one. Without it, no lead and no allocated SR can be attributed.}

\builtfrom{Screen: the Profile page's admin tab. API \texttt{GET /pms/se-uid},
\texttt{POST /pms/se-uid}, \texttt{POST /pms/se-uid/sync},
\texttt{POST /pms/se-uid/import}, \texttt{DELETE /pms/se-uid/\{id\}}. Server
\file{pms\_controller.py} -- \texttt{se\_uid\_payload}, \texttt{sync\_se\_uids},
\texttt{import\_se\_uids}. Table \file{pms\_se\_uid\_master}.}

\devproc
\begin{itemize}
  \item The Response Time \& MaxTTR file carries only the engineer's \col{se name}, while
  the LMS and EFSR files identify them by \col{service engineer uid}. The master stores
  both, keyed on the squashed upper-case name, so spacing or case differences between
  files never split one engineer into two.
  \item \textbf{Sync} scans the uploaded files for engineer names and adds the ones the
  master does not yet have, stamping which file each was found in (MaxTTR / LMS / EFSR).
  A UID can also be imported in bulk from an Excel sheet.
  \item One master row may carry more than one UID -- an engineer who appears under two
  ids in the source systems is still one person on the reports.
  \item The row also holds an optional \textbf{branch}. It is the last-resort answer for
  an engineer whose files give no valid KALA branch code, and is the only manual way to
  place such a person on a report.
  \item This master stays \emph{Master Admin only}, whatever PMS rights a user holds.
\end{itemize}

\handle{After a new import, press Sync, fill in the missing UIDs, and pin a branch for any
engineer the reports show under Unmapped Branch.}

\screenshot{screens/profile-se-uid-master.png}{Profile -- SE UID Master}

% ---------------------------------------------------------------------------
\subsection{Sales and Labour Report -- Upload and Preview}
\covers{R1, R2, R3}

\purpose{Get the Part Sale and Labour Revenue files into the system safely, and let the
user see exactly what is stored.}

\builtfrom{Screen \file{client/src/pages/SalesLabourReport.jsx} (the two upload boxes and
the \emph{Uploaded File Preview}). API \texttt{POST /pms/upload},
\texttt{GET /pms/upload/progress}, \texttt{GET /pms/uploads},
\texttt{GET /pms/data/summary}, \texttt{GET /pms/data/preview},
\texttt{POST /pms/data/cancel-row}, \texttt{DELETE /pms/data}. Server
\file{pms\_controller.py} -- \texttt{validate\_file}, \texttt{import\_file},
\texttt{preview\_rows}, \texttt{set\_row\_cancelled}. Tables
\file{pms\_sales\_records}, \file{pms\_upload\_batches}.}

\devproc
\begin{itemize}
  \item \textbf{Check file format} runs a dry run: it reports which canonical columns were
  found, which are missing and a small sample, and writes nothing. The panel lists the
  full standard column list of each file, with the three critical columns marked.
  \item The import runs in a worker thread and reports live progress, so a large file
  never blocks the page.
  \item Identity is \col{claim invoice number} plus the file's line discriminator. The
  Part Sale file has one row per part line (\col{part number}); the Labour file has one
  row per labour line -- a single SR is billed as several lines sharing one
  \col{sr number}, so the SR alone is not unique and the discriminator adds
  \col{labour part num}. Keying on the SR alone once kept only the last line of each SR
  and made the labour report read under the file's own total.
  \item Every upload returns inserted / updated / duplicate / skipped counts and is kept
  as a batch row, so the history of what was loaded is auditable.
  \item The preview shows the stored rows in the same column order as the standard Excel
  file, with a search box, and lets an invoice be \textbf{cancelled}: the row stays in the
  database and is excluded from every report until it is restored.
\end{itemize}

\handle{Check the format, upload the Part Sale and Labour files (any period, as often as
needed), read the preview to confirm, and cancel any invoice the business has withdrawn.}

\screenshot{screens/sales-upload-preview.png}{Sales and Labour -- upload boxes and the
Uploaded File Preview}

The two files are stored column for column. The layout below is also the order the
preview and the exported workbook use.

\begin{table}[H]\centering\small
\begin{tabularx}{\linewidth}{@{}p{6.5cm}X@{}}
\toprule
\textbf{Part Sale (spares) file} & \textbf{Labour Revenue file} \\
\midrule
ZONE NAME, SD ID, SD NAME, \textbf{BRANCH ID*}, BRANCH NAME, INSTANCE ID, SEGMENT,
PRODUCT SEGMENT, APPLICATION CODE, ENGINE SERIAL NO, CLAIM INVOICE NUMBER,
\textbf{CLAIM INVOICE DATE*}, CLAIM INVOICE SR TYPE, CLAIM INVOICE SR SUB TYPE, CATEGORY,
PART CATEGORY, PART NUMBER, PART DESCTRIPTION, QUANTITY,
\textbf{NET TAXABLE AMOUNT*}
&
ZONE NAME, SD ID, SD NAME, \textbf{BRANCH ID*}, BRANCH NAME, CLAIM INVOICE NUMBER,
\textbf{CLAIM INVOICE DATE*}, PRODUCT SEGMENT, SEGMENT, SERIES, SR NUMBER, SR TYPE,
SR SUBTYPE, \textbf{NET TAXABLE AMOUNT*} \\
\bottomrule
\end{tabularx}
\caption{The two sales files. * marks the three critical columns; every other column is
stored as well, and anything unrecognised is preserved rather than dropped}
\end{table}

% ---------------------------------------------------------------------------
\subsection{Sales and Labour Report -- All (Spare + Labour)}
\covers{R5, R6, R7, R8, R24}

\purpose{The headline report: every branch's Spare and Labour sale for the chosen period,
measured against the AOP, with four breakdown tables underneath.}

\builtfrom{Screen \file{SalesLabourReport.jsx} (report type \emph{All}). API
\texttt{GET /pms/report}. Server \file{pms\_controller.py} --
\texttt{generate\_report} / \texttt{\_generate\_report}. Reads
\file{pms\_sales\_records} against \file{pms\_branch\_targets},
\file{pms\_month\_settings}, \file{pms\_holidays} and \file{pms\_sr\_type\_map}.}

\devproc
\begin{itemize}
  \item The period defaults to month-to-date and can be any range -- week, month, quarter
  or year. The Part Sale and Labour sides each carry their \emph{own} as-on date, so a
  file that is a week behind is measured only against the targets up to its own last data
  date.
  \item Targets come from every month the period touches, and are prorated on
  \emph{working days} of the branch's own region -- Sundays and the ticked holidays carry
  no target.
  \item The whole aggregation is done in SQL: branch, region, segment, head and category
  buckets come back as a few dozen grouped rows rather than every sales row travelling
  over the network. Invoice counts are distinct counts of \col{claim invoice number}
  (falling back to the row id when the file left it blank).
  \item Cancelled invoices are excluded everywhere.
\end{itemize}

The branch tables (one for Spare, one for Labour) carry these columns, and every figure
is spelled out on the screen so the reader never has to guess a formula:

\begin{table}[H]\centering\small
\begin{tabularx}{\linewidth}{@{}p{4.4cm}X@{}}
\toprule
\textbf{Column} & \textbf{How it is worked out} \\
\midrule
Responsible Person & the branch's Branch Admin, from the AOP Target Master \\
Branch ID / Branch Name / Region & the AOP Target Master, falling back to the file \\
Target (period) & the full monthly targets of every month the period touches, summed \\
Daily Target & Target $\div$ the working days of the period, in the branch's own region \\
Achi. on \textit{date} & \col{net taxable amount} on the period's last day \\
Target till \textit{date} & each month's target $\times$ (its working days inside the
period $\div$ its working days) \\
Achi. till \textit{date} & \col{net taxable amount} summed over the period \\
Invoice Count till & distinct \col{claim invoice number} in the period \\
\% Achieved & Achi. till $\div$ Target till $\times$ 100 \\
Short-Fall & Target till $-$ Achi. till (a negative figure means ahead of plan) \\
Balance / Month & Target $-$ Achi. till \\
\bottomrule
\end{tabularx}
\caption{Sales and Labour -- the branch table columns}
\end{table}

Four breakdown tables follow, all for the same period:

\begin{itemize}
  \item \textbf{Region-wise Sales} -- MH first, then KA, in a fixed order, each with its
  own AOP target in front of the sale, because a region target is simply the sum of its
  branches' targets.
  \item \textbf{Segment-wise Sales} -- the row's \col{segment}, falling back to
  \col{product segment}; a row with neither is OTC business.
  \item \textbf{Service Report Type-wise Sales} -- the head an SR Type maps to in the SR
  Type Master (Sales and Labour). Segment and Service Head have no target dimension in the
  AOP, so their \% column is a share of the column total rather than an achievement.
  \item \textbf{Category-wise Sales (Spare)} -- spares only, since the Labour file has no
  \col{category} column: sale, quantity, line items and invoice count per category.
\end{itemize}

\handle{Set the period, press Generate, read the two branch tables, then the breakdowns.
Export to Excel for a copy that mirrors the screen -- Spare and Labour side by side,
region by region.}

\screenshot{screens/sales-report-all.png}{Sales and Labour -- All (Spare + Labour): the
branch tables and the breakdowns}

% ---------------------------------------------------------------------------
\subsection{Sales and Labour Report -- Branch-wise Report}
\covers{R9, R24}

\purpose{The same period, but for a hand-picked set of branches, with the detail a branch
review actually needs.}

\builtfrom{Screen \file{SalesLabourReport.jsx} (report type \emph{Branch-wise}). API
\texttt{GET /pms/report/branch-detail}. Server \texttt{branch\_detail\_report} /
\texttt{\_branch\_detail\_report}.}

\devproc
\begin{itemize}
  \item The view opens with every branch pre-selected and narrows as branches are ticked;
  the branch tables above shrink to the selection.
  \item Five detail sections are recomputed from \emph{only} those branches' rows:
  \textbf{Week-wise} (the period cut into seven-day chunks, with Spare, Labour and invoice
  count), \textbf{Month-wise}, \textbf{Segment-wise}, \textbf{Service Report Type-wise} and
  \textbf{Category-wise (Spare)}.
  \item Each week also carries its slice of the AOP target: a month's target spread over
  its working days and the days of that week added up -- the same calendar the main report
  uses, so Sundays and ticked holidays carry no target. When the selection spans both
  regions a day counts as off only when it is off in every region selected.
\end{itemize}

\handle{Choose Branch-wise, tick the branches, open the sections needed, export.}

\screenshot{screens/sales-report-branchwise.png}{Sales and Labour -- Branch-wise Report}

% ---------------------------------------------------------------------------
\subsection{Sales and Labour Report -- FY / Quarterly / Month-wise}
\covers{R10, R24}

\purpose{The whole financial year on one screen: how the year is tracking, quarter by
quarter and month by month, and what run-rate is needed to close it.}

\builtfrom{Screen \file{SalesLabourReport.jsx} (report type \emph{Spares and Labour (FY /
Quarterly / Month)}). API \texttt{GET /pms/report/fy-summary}. Server
\texttt{fy\_summary\_report} / \texttt{\_fy\_summary\_report}.}

\devproc
\begin{itemize}
  \item These three panels ignore the period picker entirely -- they always cover the full
  financial year that is asked for, April to March.
  \item Targets are the AOP Master's \emph{full} monthly values, never prorated here;
  achievement is aggregated in SQL by branch, month and record type, with cancelled rows
  excluded.
  \item The report also returns how far into the year the sales data actually reaches, and
  the year's working days split into elapsed and remaining, per region. That is what makes
  the run-rate columns honest: they rest on the same calendar the targets do.
  \item Quarters are the financial ones -- Q1 is April to June, Q4 is January to March.
  \item The export writes three sheets in one workbook: the FY table, Quarterly and
  Month-wise, with Spare and Labour across the year.
\end{itemize}

\handle{Pick the financial year, open the FY, Quarterly or Month panel, export.}

\screenshot{screens/sales-report-fy.png}{Sales and Labour -- FY / Quarterly / Month-wise}

% ---------------------------------------------------------------------------
\subsection{Employee Productivity (Service Engineer Productivity)}
\covers{R12, R13, R14, R15, R24, R27}

\purpose{One row per service engineer answering: how much did they close, how many days
did they work, how happy were the customers, and how much business did they raise.}

\builtfrom{Page \file{client/src/pages/EmployeeProductivity.jsx} (owns the period picker);
table \file{client/src/components/EmployeeProductivityReport.jsx}. API
\texttt{GET /pms/report/employee-productivity}. Server \file{pms\_controller.py} --
\texttt{employee\_productivity\_data} / \texttt{\_employee\_productivity\_data}.}

\textbf{The MaxTTR file is the source.} Every SR in the \emph{Response Time \& MaxTTR
Details} import that has a close date and an engineer name counts once for that engineer
on that date. The labour file is not involved: an SR is counted because the engineer
closed it, not because it was billed. Four files and four masters then hang off that
spine.

\begin{table}[H]\centering\small
\begin{tabularx}{\linewidth}{@{}p{3.1cm}p{4.3cm}X@{}}
\toprule
\textbf{Group of columns} & \textbf{Source file} & \textbf{Columns read} \\
\midrule
Close SR, and its SR Type split &
Response Time \& MaxTTR Details &
\col{se name}, \col{branch id}, \col{branch name}, \col{sr close date}, \col{sr type}
(through the SR Type Master (MaxTTR)) \\
\addlinespace
Days present on Task end &
Response Time \& MaxTTR Details &
distinct \col{sr task end date} per engineer \\
\addlinespace
Allocated SR and its split &
EFSR Report &
\col{service engineer uid}, \col{sd branch code}, \col{task assigned date},
\col{sr type} (through the SR Type Master (EFSR)) \\
\addlinespace
CDI (Passive, Detractor, \%) &
CDI Detail Report &
\col{x technician name}, \col{activity end date}, \col{cdi category} \\
\addlinespace
Number of Lead, Product Wise Lead Count, Spare / Labour Conv. Amount &
LMS Data for ERP &
\col{lead number}, \col{service engineer uid}, \col{branch id},
\col{lead created date}, \col{lead raised for} (through the Lead Category Master),
\col{part invoice amount}, \col{labour invoice amount} \\
\addlinespace
Working Days &
AOP and Master &
the month-wise working-days master, per region, prorated to the picked period \\
\bottomrule
\end{tabularx}
\caption{Employee Productivity -- every column and the column it comes from}
\end{table}

\devproc
\begin{itemize}
  \item \textbf{Identity.} An engineer is their letters-only name plus their branch, so
  \emph{Vijaykumar Jadhav}, \emph{VijaykumarJadhav} and \emph{VIJAYKUMAR  JADHAV} are
  always one engineer however the files spell them; the best-spaced spelling seen is kept
  for display.
  \item \textbf{Branch} is the MaxTTR row's own \col{branch id}, never the branch name --
  files spell the same branch differently and the spelling changes between imports, while
  the id does not. The name is only a label, and the AOP master's spelling is used so this
  report and the Sales report label a branch identically.
  \item \textbf{Leads} are attributed strictly through the SE UID Master: an LMS row's
  \col{service engineer uid} is looked up there to find the engineer, and is never matched
  against the MaxTTR names. A UID in no master row is counted as an \emph{unlinked lead}
  and reported, so the master can be fixed. Leads are counted once per distinct
  \col{lead number}.
  \item \textbf{Allocated SR} is counted on \col{task assigned date}, not on a closed
  date: the closed date is null for exactly the SRs this column is meant to show -- the
  ones assigned but not yet finished. (The columns are hidden on the current screen; the
  payload and every rollup still carry them.)
  \item \textbf{Days present} counts distinct \col{sr task end date} values -- the day the
  engineer finished the job in the field. The close date is an office event that often
  lands days later and in a batch, so it would overstate attendance.
  \item \textbf{Productivity} is Close SR $\div$ \emph{Working Days}, not $\div$ days
  present: days present is attendance, while the target is the available man-days from the
  AOP master. Hiding the Working Days column does not change the division.
  \item \textbf{CDI \%} is (Promotor $-$ Detractor) $\div$ all feedback $\times$ 100. The
  screen shows Passive and Detractor; Promotor is still counted and still drives the \%.
  \item \textbf{Working Days and Days present print on engineer rows only.} Every rollup
  row -- branch, group, region, Grand Total -- leaves both blank rather than summing
  man-days into a number nobody would read.
  \item \textbf{Rollups.} Branches roll into the configured branch groups, then into MH /
  KA, then into the Grand Total. A branch code no KALA master knows folds into the single
  \qt{Unmapped Branch} bucket, so the report still adds up to the file.
  \item \textbf{One fetch.} The payload is raw per-day records; changing the period, the
  weekly split or any filter re-aggregates in the browser without another request.
\end{itemize}

\handle{Set the period at the top, use the toolbar to show the weekly split, hide Working
Days, or fold the SR Type / CDI / Lead Count bands, then export. The \% and count columns
carry their formula in the header, so a figure can always be traced.}

\screenshot{screens/employee-productivity.png}{Employee Productivity -- service engineer
rows with SR Type, CDI and lead bands}

% ---------------------------------------------------------------------------
\subsection{SR Allocation Report}
\covers{R16, R24, R27}

\purpose{What was given to each engineer and what they finished -- and therefore what is
still open.}

\builtfrom{Page \file{client/src/pages/SRAllocation.jsx}; table
\file{client/src/components/SRAllocationReport.jsx}. API
\texttt{GET /pms/report/sr-allocation}. Server \texttt{sr\_allocation\_data} /
\texttt{\_sr\_allocation\_data}.}

\begin{table}[H]\centering\small
\begin{tabularx}{\linewidth}{@{}p{3.4cm}p{3.6cm}X@{}}
\toprule
\textbf{Measure} & \textbf{Source file} & \textbf{Columns read} \\
\midrule
Allocated Total Task & EFSR Report &
\col{service engineer uid}, \col{task assigned date}, \col{sr type},
\col{sd branch code} \\
Closed Total Task & EFSR Report &
the same rows, counted on \col{task end date} \\
Open SR (branch wise) & Open SR Load Report &
\col{instance id [asset \#]} joined to the customer master for its \col{branch id}, and
\col{sr created date} \\
Home branch fallback & Response Time \& MaxTTR &
\col{se name}, \col{branch id} \\
SE Name / SE UID & SE UID Master & UID to the engineer's canonical name \\
SR Type rows & SR Type Master (EFSR) & the head each \col{sr type} maps to \\
\bottomrule
\end{tabularx}
\caption{SR Allocation -- sources and columns}
\end{table}

\devproc
\begin{itemize}
  \item One EFSR row normally counts \emph{twice}: once on its assigned date and once on
  its task end date. The two totals differ by the SRs assigned inside the period but
  closed outside it, or not closed at all -- and that gap is the open backlog.
  \item \textbf{Closed} deliberately uses \col{task end date}, the day the engineer
  finished, not the back-office \col{sr closed date}, which lands days later and never
  arrives at all for about a third of the tasks the engineer has actually completed.
  \item Rows are attributed through the SE UID Master, so the SE Name and SE UID read
  exactly the same here as on the Employee Productivity report. A UID the master does not
  know still counts, under the file's own engineer name, so the report always adds up to
  the file; how many such rows there are is reported.
  \item \textbf{Home branch.} The EFSR file carries the SD branch code of the SR, and on
  about 1\% of rows that is another dealer's code. Grouping straight on it would scatter
  one engineer across ghost branches, so the report groups by the \emph{engineer}: every
  SR of theirs counts in the branch they belong to, resolved as (1) the KALA branch they
  have the most EFSR rows in, (2) the branch the MaxTTR file puts them in, (3) the branch
  pinned on their SE UID Master row, (4) the single Unmapped Branch bucket.
  \item The \textbf{Open SR} grid uses exactly two columns of the Open SR Load Report and
  counts the whole accumulated table -- it answers \qt{SRs open in this branch}, not
  \qt{in today's upload}.
  \item The date columns can be bucketed by day, week or month, and the \emph{Date split}
  control chooses whether they carry Allocated, Closed or both; the two totals are always
  shown.
\end{itemize}

\handle{Set the period, choose the granularity and the date split, open an engineer's
row for the SR-Type-wise breakdown, export.}

\screenshot{screens/sr-allocation.png}{SR Allocation Report -- allocated vs closed, engineer
and date wise}

% ---------------------------------------------------------------------------
\subsection{Training Report}
\covers{R17, R24}

\purpose{The service engineers' skill and training master, readable from either end:
\qt{what has this person been trained on} and \qt{who holds this skill}.}

\builtfrom{Screen \file{client/src/pages/TrainingReport.jsx}. API
\texttt{POST /pms/training/upload}, \texttt{GET /pms/training/report},
\texttt{GET /pms/training/summary}, \texttt{DELETE /pms/training/data}. Server
\file{server/app/controllers/pms\_training\_controller.py}. Table
\file{pms\_training\_records}; the upload audit shares \file{pms\_upload\_batches} with
\texttt{record\_type = training}.}

\devproc
\begin{itemize}
  \item \textbf{Eight columns are fixed} and stored as real columns: \col{uid no},
  \col{employee ticket number}, \col{full name}, \col{occupation}, \col{skill},
  \col{branch name}, \col{branch id}, \col{hire date}. Only \col{uid no} and
  \col{full name} are critical.
  \item \textbf{Every other column of the file is dynamic} -- zone name, job title,
  category, training date, training end date, SD name / id, the bank block and so on --
  kept verbatim and rendered from there, so the export can gain or lose a column without
  a code change. The server then splits the dynamic columns into the ones that describe
  the \emph{person} and the ones that describe the \emph{training}.
  \item \textbf{Identity} is \col{uid no} + \col{skill} + \col{training date} +
  \col{category}. Verified against the real 407-row export: UID + skill alone collapsed
  114 rows, and adding the training date still collapsed 107, because the same training is
  recorded under more than one category; all four together drop none. The two extra
  columns are read for the key even though they stay dynamic -- the key has to be stable,
  the storage does not.
  \item It is a \emph{master}, not a period report: rows accumulate across uploads and a
  re-uploaded row with changed values updates in place, so the page always shows the whole
  training history. An engineer with no training yet still gets a row, with the skill
  blank, so the master lists everyone the file carries.
  \item \textbf{By Employee} lists one row per person; clicking opens their skills first,
  then every other detail the file carries. \textbf{By Skill} pivots the same payload in
  the browser and lists one row per skill; clicking opens everyone who holds it. Searching
  either way is instant and never hits the server.
  \item The search compares on the squashed name form, so \qt{nilesh salunke},
  \qt{NILESHSALUNKE} and \qt{Nilesh  Salunke} all find the same person.
\end{itemize}

\handle{Upload the Training Report export whenever KOEL issues a new one, then search by
employee name or by skill and click the row for the full record. Export to Excel as
needed.}

\screenshot{screens/training-report.png}{Training Report -- By Employee and By Skill}

% ---------------------------------------------------------------------------
\subsection{Annual Reports -- the page}
\covers{R18, R19, R20, R21, R22}

\purpose{One page holding the yearly management sheets. The report picker chooses which
sheet is shown; the period picker at the top owns the reporting period for all of them.}

\builtfrom{Page \file{client/src/pages/AnnualReports.jsx}, with one component per sheet:
\file{ServicePenetrationReport.jsx}, \file{AmcBandhanProjectionReport.jsx},
\file{AmcMonthlyReport.jsx}, \file{CustomerDelightIndexReport.jsx},
\file{ServiceLoadResponseReport.jsx}; shared workbook chrome in
\file{client/src/utils/pmsExport.js} and shared table chrome in
\file{client/src/components/reportChrome.jsx}.}

Each sheet has its own source file, so each brings its own date range: the period picker
re-defaults to the financial year of whichever report is open, and offers only the
financial years that report's file actually covers. The AMC report is two tabs of one
report rather than two entries in the picker, because both tabs read the same file.

% ---------------------------------------------------------------------------
\subsection{Annual Reports -- Service Penetration}
\covers{R18}

\purpose{How much of the installed base was actually touched by a closed service call.}

\builtfrom{\file{ServicePenetrationReport.jsx}; API
\texttt{GET /pms/report/annual/service-penetration}; server
\texttt{\_annual\_service\_penetration\_data}.}

\begin{table}[H]\centering\small
\begin{tabularx}{\linewidth}{@{}p{3.3cm}p{3.4cm}X@{}}
\toprule
\textbf{Block} & \textbf{Source file} & \textbf{Columns read} \\
\midrule
Ind / PG Population, Grand Total & Asset Detailed Report &
\col{branch id}, \col{segment} (IND / PG), \col{commissioning date},
\col{asset operational status} \\
UNQ AST SR CLOSED & Response Time \& MaxTTR &
\col{branch id}, \col{segment}, \col{instance id} (or \col{engine serial no}),
\col{sr close date} \\
Pen \% & derived & closed $\div$ population $\times$ 100 \\
\bottomrule
\end{tabularx}
\caption{Service Penetration -- sources and columns}
\end{table}

\devproc
\begin{itemize}
  \item The population is \emph{cumulative}: each asset counts once, in its own branch, on
  its commissioning date, and the sheet reads the installed base \emph{as on} the end of
  the period, not what was commissioned inside it.
  \item Assets the ERP has retired (operational status \emph{Inactive}) are out of the
  base; an asset with no commissioning date cannot be placed on the calendar and is
  reported rather than dropped.
  \item \textbf{Unique means unique}: an asset with five closed SRs counts once. Because a
  distinct-asset count cannot be added up day by day the way a population can, the payload
  ships one row per asset carrying the days it was serviced on, and the report counts the
  assets with at least one day inside the period. Every period is therefore exact and
  still needs no refetch.
  \item An asset is counted under the branch and segment its \emph{own SR row} names, so
  nothing is lost when the asset master has not caught up yet.
  \item Only the Total Pen \% cell is shaded, against a threshold that is deliberately not
  printed anywhere on the sheet -- it is a fill rule, not a figure.
\end{itemize}

\handle{Pick the period; read population, serviced assets and Pen \% per branch, with MH,
KA and overall totals. Export mirrors the screen.}

\screenshot{screens/annual-service-penetration.png}{Annual Reports -- Service Penetration}

% ---------------------------------------------------------------------------
\subsection{Annual Reports -- AMC and Bandhan Projection (tab 1)}
\covers{R19, R23}

\purpose{Per branch: the year's AMC position against the plan, and the month's business
actually done.}

\builtfrom{\file{AmcBandhanProjectionReport.jsx}; API
\texttt{GET /pms/report/annual/amc-bandhan}; server \texttt{\_annual\_amc\_bandhan\_data}.}

\begin{table}[H]\centering\small
\begin{tabularx}{\linewidth}{@{}p{4.6cm}p{3.4cm}X@{}}
\toprule
\textbf{Column} & \textbf{Source} & \textbf{How it is produced} \\
\midrule
F\textit{nn} ACT D/BAMC & AOP and Master & typed -- last year's actual \\
F\textit{nn} PROJ AOP D/BAMC & AOP and Master & typed -- this year's projection \\
BEST ACT AOP D/BAMC (M) & AOP and Master & typed -- the year's best month \\
F\textit{nn} YTD ACT AMC NOS & AMC Population Report &
counted: \col{agreement status} is Active \emph{and} \col{agreement end date} is not
past -- a population, so neither the year nor the month moves it \\
\textit{Mon}-\textit{yy} Act & the four Bandhan quote files &
counted: one row per quote on \col{payment update date and time}, branch wise \\
\bottomrule
\end{tabularx}
\caption{AMC and Bandhan Projection -- where each column comes from}
\end{table}

\devproc
\begin{itemize}
  \item The sheet asks two different questions and no one file answers both, so it reads
  two sources: the agreement population and the money received.
  \item \textbf{The month's business is payments, not cover starting.} Counted on the
  quote files' payment date the figures tie up with the business's own sheet (Mar-26 145,
  Apr-26 105, Jun-26 104 against 144 / 108 / 106), whereas the agreements' start months
  are 15 to 30 out. This is the only column the month picker moves.
  \item \textbf{The first three columns are the AOP master's and nothing else's.} Last
  year's actual could not be counted honestly anyway: the AMC file is a snapshot keyed on
  instance id, one row per genset carrying only its latest agreement, so a renewal
  overwrites the agreement it replaced -- FY26 counts 937 from today's file against the
  1,353 the business closed it with.
  \item \textbf{Best month} follows the monthly payment series, and a region or company
  row takes the best month of \emph{its own} monthly totals -- never the sum of its
  branches' best months, which fall in different months and would invent a month nobody
  had.
  \item Every other rollup is a plain sum of the branches. The typed figures stay null
  until at least one branch has one, so an unfilled region reads as empty rather than as a
  target of zero.
\end{itemize}

\handle{Open the AMC report, stay on the Projection tab, pick how far into the year to
read, export.}

\screenshot{screens/annual-amc-projection.png}{Annual Reports -- AMC and Bandhan Projection}

% ---------------------------------------------------------------------------
\subsection{Annual Reports -- AMC (tab 2, the monthly sheet)}
\covers{R20, R23}

\purpose{The same AMC file read by agreement \emph{category} instead of by branch, one
financial year wide.}

\builtfrom{\file{AmcMonthlyReport.jsx}; API \texttt{GET /pms/report/annual/amc-monthly};
server \texttt{\_annual\_amc\_monthly\_data}. Reads \col{branch id},
\col{agreement type}, \col{agreement start date}, \col{agreement end date},
\col{agreement status} and \col{last agreement number} from the AMC Population Report.}

\begin{table}[H]\centering\small
\begin{tabularx}{\linewidth}{@{}p{5.4cm}X@{}}
\toprule
\textbf{Row} & \textbf{What it counts} \\
\midrule
KOEL Bandhan Total & the MH and KAR rows added up \\
KOEL Bandhan MH / KAR & the KOEL Bandhan family, split on the branch's region \\
KALA AMC & agreement type \emph{Dealer Agreement} -- the dealership's own AMC \\
KOEL Corporate AMC & agreement type \emph{KOEL Agreement}, printed cumulative \\
AMC Expired During the Month & counted on \col{agreement end date} \\
AMC Renewed During the Month & rows that carry a \col{last agreement number} \\
Live AMC (KOEL+KALA) OVERALL & KOEL Bandhan + KALA AMC \\
\bottomrule
\end{tabularx}
\caption{The AMC monthly sheet's rows}
\end{table}

\devproc
\begin{itemize}
  \item Columns are: the AOP (from AOP and Master), the year so far, one column per month
  reached, and the working month \emph{also} broken into its Monday-to-Sunday weeks --
  beside its month column, not instead of it, which is how the business prints the sheet.
  \item \textbf{Closed years are deliberately not shown.} Counted from today's snapshot
  the Bandhan rows give 235 for FY24--25 and 858 for FY25--26 against the 837 and 1,302
  the business closed those years with. A column that wrong is worse than no column, so
  the sheet prints only the running year, which is counted live and is sound.
  \item Any agreement type the row mapping does not know is reported rather than dropped,
  so a new type from KOEL shows up instead of vanishing.
  \item The second table, \emph{KCGL Total AMC}, plans against the sum of every branch's
  projection from the AOP master's first table, spread equally across the twelve months --
  so the yearly plan lives in exactly one place and the company row is derived from it
  rather than typed a second time and left to drift.
\end{itemize}

\handle{Open the AMC report and switch to the AMC tab. The AOP comes from AOP and
Master $\rightarrow$ AMC and Bandhan AOP (the category table).}

\screenshot{screens/annual-amc-monthly.png}{Annual Reports -- the AMC monthly sheet}

% ---------------------------------------------------------------------------
\subsection{Annual Reports -- Customer Delight Index (CDI)}
\covers{R21, R23}

\purpose{The customer feedback score per branch, against its AOP, for the financial year.}

\builtfrom{\file{CustomerDelightIndexReport.jsx}; API \texttt{GET /pms/report/annual/cdi};
server \texttt{\_annual\_cdi\_data}. Reads \col{branch name}, \col{activity end date} and
\col{cdi category} from the CDI Detail Report, and the AOP from
\file{pms\_cdi\_targets}.}

\devproc
\begin{itemize}
  \item Every feedback row falls into one of three buckets, read off the \emph{leading
  word} of \col{cdi category} so the score range in brackets can change without breaking:
  \emph{Promotor(09-10)}, \emph{Detractor(00 - 06)}, and everything else Passive.
  \item The score of any set of rows is (Promotor $-$ Detractor) $\div$ (Promotor +
  Passive + Detractor) $\times$ 100. It is therefore \textbf{not} an average of the rows
  beneath it: every total row is scored from its own feedback counts, which is why a
  region row can sit outside the range of its branches.
  \item The CDI Detail Report is the only PMS import that carries no branch id -- it names
  its branch (\qt{KALA Care Global LLP - Ahmednagar}) and nothing else. The name-to-id map
  is therefore built \emph{from the other files}, which carry both, so when the ERP renames
  a branch the map follows on the next upload with no code change. The AOP master's own
  spelling is registered against the same id as well, so either wording resolves.
  \item Feedback whose branch no master knows folds into the Unmapped Branch bucket, and
  rows with no activity end date cannot be placed on the calendar at all; both are
  reported.
  \item The columns follow the financial year the period lands in: the two previous
  financial years as cumulative columns, the AOP, the year so far, one column per whole
  month, and the working month split into its Monday-to-Sunday weeks. A finished year
  reads as twelve month columns and no weeks.
\end{itemize}

\handle{Pick the financial year, read branch, region and overall scores against the AOP,
export.}

\screenshot{screens/annual-cdi.png}{Annual Reports -- Customer Delight Index}

% ---------------------------------------------------------------------------
\subsection{Annual Reports -- Service Load and Response}
\covers{R22, R23, R24}

\purpose{The service department's own scorecard: how much load was closed, how fast it was
responded to, and how productive the engineers were.}

\builtfrom{\file{ServiceLoadResponseReport.jsx}; API
\texttt{GET /pms/report/annual/service-load}; server
\texttt{\_annual\_service\_load\_data}. \textbf{One file only} -- Response Time \& MaxTTR
Details -- with three masters on top: the SR Type Master (Service Load), the Service Load
AOP targets and the typed FTR / FVR figures.}

The printed sheet is four pages; the screen is \textbf{six tabs}, because its response page
carries two independent measures and its closing blocks summarise the whole sheet:

\begin{table}[H]\centering\small
\begin{tabularx}{\linewidth}{@{}p{3.9cm}X@{}}
\toprule
\textbf{Tab} & \textbf{What it shows} \\
\midrule
Service Load & the SR-type-wise breakdown and its Total, then the same load by branch
under MH / KA, closing on the Grand Total \\
Productivity & Calls PP PD per region, branch and overall, plus the
\emph{Service Request Closure (Nos.)} month strip \\
4 Hrs Response & 4 HRS RESPONSE \% per region and branch \\
SR Closed 24 Hrs & SR CLOSED WITHIN 24 HRS \% \\
SR Closed 48 Hrs & SR CLOSED WITHIN 48 HRS \% \\
Overall and FTR/FVR & MAX TTR Overall, and the FTR / FVR figures typed in AOP and Master \\
\bottomrule
\end{tabularx}
\caption{Service Load and Response -- the six tabs}
\end{table}

\devproc
\begin{itemize}
  \item Everything is counted on \col{sr close date}, per \col{branch id}. \emph{Total
  Service Load Available} is every SR of the period whatever its type; the breakdown rows
  are that same population split by the head master. An SR Type with no head still counts
  in the total but gets no breakdown row, so the gap is visible on the sheet rather than
  swallowed.
  \item \textbf{All three compliance percentages read one column} -- \col{response time} --
  over every SR of the period: $\le 4$, $\le 24$, $\le 48$ hours.
  \item \textbf{\col{sr open date} is never used}, and this is the single most important
  decision on the sheet. In the live file 8,721 of 45,279 rows carry an open date
  \emph{after} their close date, always at midnight: they are the scheduled dates of
  planned work that was closed early. The file's own arithmetic turns those into a
  negative response time. So response and TTR are read from the file's own columns and
  never recomputed from open to close.
  \item \textbf{A negative response time counts as zero hours.} The job was closed before
  it was due, so it is inside every window -- which is exactly how the source file scores
  it. Clamping rather than dropping keeps the denominator equal to the load count, so the
  sheet's SR total and its percentages describe one population; a cell says on hover how
  many of its SRs were clamped.
  \item \textbf{Every percentage row is scored from its own SRs}, never averaged from the
  rows beneath it -- a region row divides the region's compliant SRs by the region's SRs.
  Productivity is the same: a region's closures over the region's man-days.
  \item \textbf{Productivity} is closures $\div$ (working days $\times$ the engineers who
  closed something in that window). A distinct engineer count cannot be summed across
  days, so the payload ships one triple per engineer-day and the page counts the distinct
  engineers inside each window itself.
  \item \textbf{FTR / FVR are not counted.} First Time Right and First Visit Report are
  judgements KOEL makes about a job, and the two files that once carried them were
  withdrawn -- so the sheet prints what a person entered in AOP and Master, and a period
  nobody filled in shows a dash.
  \item 45k rows are aggregated in SQL down to (branch, day, SR type) -- about 12k records
  -- and the FY cumulatives, month columns, weekly split and every region rollup are
  derived in the browser from those, with no second request.
\end{itemize}

\handle{Pick the financial year, switch tabs for load, productivity and each response
window, and read the AOP column beside every measure. Export writes the tab you are on.}

\screenshot{screens/annual-service-load.png}{Annual Reports -- Service Load and Response}

% ============================== 7. CROSS-CUTTING ============================
\section{Cross-cutting Behaviour}

\subsection{Access rights}
\covers{R25}
PMS access is a per-user flag; the AOP and Master page is granted per tab at
\emph{none} / \emph{view} / \emph{edit}. On a view-only tab the save, sync, reset and
master buttons are hidden and every cell is locked -- and the server enforces the same per
tab, so a user granted only one tab cannot read or write another tab's data by calling its
endpoint directly. Three registries hold the tab keys and must stay in step:
\file{pagePermission.js} on the client, \file{user\_controller.py} and the \texttt{tab=}
argument of \texttt{\_require\_aop} in \file{pms\_routes.py}. The SE UID Master stays
Master-Admin-only whatever PMS rights a user holds.

\subsection{Performance on a remote database}
\covers{R26}
The database sits across the network, so the rule is: aggregate in SQL, ship the small
result. Every report groups its rows server-side and returns raw per-day records, and the
page derives periods, granularities, weeks and rollups from those without another request.
Results are cached against a fingerprint of the data; an upload changes the fingerprint
and the next report rebuilds. A new PMS report must add its tables to that fingerprint
query, or it will serve a stale cache after an import.

\subsection{Nothing is dropped silently}
Every report reports its own gaps rather than hiding them: unmapped SR Types and lead
values get a visible \qt{Unmapped} bucket, branch codes no KALA master knows fold into one
\qt{Unmapped Branch} row, LMS UIDs missing from the SE UID Master are counted as unlinked
leads, rows with no usable date are reported as such, and clamped or excluded rows are
named in the report's meta. A figure the module cannot produce honestly is typed by a
person and marked as typed.

\subsection{Excel export}
\covers{R24}
Every report exports through the shared chrome in \file{client/src/utils/pmsExport.js}: one
ERP blue band carrying the report and its period, a grey-blue header row, branch rows, the
MH / KA region totals and a dark overall row. The workbook uses the light palette as
literal colours -- it is read and printed on its own and must not follow the app's dark
theme. \file{exceljs} is loaded on demand, so opening a page never pays for it. Date
columns are written as real Excel dates through the shared date utility, because a date
written the naive way lands ten seconds off midnight and silently breaks Excel's date
filters. Export respects the \file{can\_export} permission.

\subsection{The colour language of the reports}
The PMS grids share one palette, defined once as CSS variables so screen and workbook
agree and both themes work:

\begin{table}[H]\centering\small
\begin{tabularx}{\linewidth}{@{}p{3.4cm}p{2.4cm}X@{}}
\toprule
\textbf{Element} & \textbf{Light value} & \textbf{Meaning} \\
\midrule
ERP band, Grand Total & \texttt{\#2F3192} / \texttt{\#23255F} & the report's own identity
row \\
Header cells & \texttt{\#E8F3FC} & every column header \\
Row A / Row B & \texttt{\#FBFDFF} / \texttt{\#F1F8FE} & engineer or branch rows \\
Sub Total / Group Total & \texttt{\#DCEBF9} / \texttt{\#CBE1F5} & a group heading its own
rows \\
Region total (MH / KA) & \texttt{\#5E8FC2} & region rollups, white text \\
SR Type block & \texttt{\#E4F0FB} & the type-wise band \\
Selected row & \texttt{\#BFDCF7} & the open or picked row \\
\% at or above target & \texttt{\#86EFAC} & green \\
\% within a fifth of target & \texttt{\#FDE047} & yellow \\
\% further below & \texttt{\#FDBA74} & amber \\
\bottomrule
\end{tabularx}
\caption{The shared PMS colour language}
\end{table}

A dash is never painted -- an absent figure stays neutral, because \qt{not entered} is not
a performance judgement.

\subsection{India-local timestamps}
\covers{R28}
Every stored timestamp is India local time, written through the module's own time helper,
so meeting dates, upload stamps and audit trails read the same for every user whatever the
server or browser timezone.

% ============================== 8. FILE MAP =================================
\section{File Map -- Which File is Responsible for What}

\begin{longtable}{@{}p{4.2cm}p{4.8cm}p{5.8cm}@{}}
\caption{PMS pages and the files that implement them}\\
\toprule
\textbf{Page / Report} & \textbf{Frontend} & \textbf{Backend} \\
\midrule
\endfirsthead
\multicolumn{3}{@{}l}{\small\itshape Table \thetable{} -- continued}\\
\toprule
\textbf{Page / Report} & \textbf{Frontend} & \textbf{Backend} \\
\midrule
\endhead
\bottomrule
\endfoot
AOP and Master\newline\file{/aop-master} &
\file{pages/AOPMaster.jsx},\newline \file{utils/pagePermission.js} &
\file{routes/pms\_routes.py},\newline \file{controllers/pms\_controller.py},\newline
\file{models/pms\_model.py} \\
\addlinespace
Sales and Labour Report\newline\file{/sales-labour-report} &
\file{pages/SalesLabourReport.jsx} &
\texttt{import\_file}, \texttt{\_generate\_report}, \texttt{\_fy\_summary\_report},
\texttt{\_branch\_detail\_report} \\
\addlinespace
Employee Productivity\newline\file{/employee-productivity} &
\file{pages/EmployeeProductivity.jsx},\newline
\file{components/EmployeeProductivityReport.jsx} &
\texttt{\_employee\_productivity\_data} \\
\addlinespace
SR Allocation Report\newline\file{/sr-allocation} &
\file{pages/SRAllocation.jsx},\newline \file{components/SRAllocationReport.jsx} &
\texttt{\_sr\_allocation\_data} \\
\addlinespace
Training Report\newline\file{/training-report} &
\file{pages/TrainingReport.jsx} &
\file{controllers/pms\_training\_controller.py} \\
\addlinespace
Annual: Service Penetration &
\file{components/ServicePenetrationReport.jsx} &
\texttt{\_annual\_service\_penetration\_data} \\
\addlinespace
Annual: AMC and Bandhan &
\file{components/AmcBandhanProjectionReport.jsx} &
\texttt{\_annual\_amc\_bandhan\_data} \\
\addlinespace
Annual: AMC monthly &
\file{components/AmcMonthlyReport.jsx} &
\texttt{\_annual\_amc\_monthly\_data} \\
\addlinespace
Annual: CDI &
\file{components/CustomerDelightIndexReport.jsx} &
\texttt{\_annual\_cdi\_data} \\
\addlinespace
Annual: Service Load &
\file{components/ServiceLoadResponseReport.jsx} &
\texttt{\_annual\_service\_load\_data} \\
\addlinespace
Shared chrome / export &
\file{components/reportChrome.jsx},\newline \file{utils/pmsExport.js},\newline
\file{utils/excelDateColumns.js} & --- \\
\addlinespace
Routing and navigation &
\file{App.jsx}, \file{components/Navbar.jsx} & --- \\
\end{longtable}

% ============================== APPENDIX ====================================
\appendix
\section{API Endpoint Map}

Every endpoint below is mounted under \file{/pms} and defined in
\file{server/app/routes/pms\_routes.py}. All of them require PMS access; the AOP endpoints
additionally check the tab they belong to.

\begin{longtable}{@{}p{7.4cm}p{8.0cm}@{}}
\caption{PMS API endpoints}\\
\toprule
\textbf{Endpoint} & \textbf{Serves} \\
\midrule
\endfirsthead
\multicolumn{2}{@{}l}{\small\itshape Table \thetable{} -- continued}\\
\toprule
\textbf{Endpoint} & \textbf{Serves} \\
\midrule
\endhead
\bottomrule
\endfoot
\texttt{GET/POST /targets/year}, \texttt{/targets/year/bulk} & Target Master grid \\
\texttt{GET/POST /holidays} & the working-days calendar \\
\texttt{/sr-types}, \texttt{/heads} (+ \texttt{/sync}, \texttt{/reset}) & SR Type Master
(Sales and Labour) \\
\texttt{/maxttr-sr-types}, \texttt{/maxttr-heads} & SR Type Master (MaxTTR) \\
\texttt{/efsr-sr-types}, \texttt{/efsr-heads} & SR Type Master (EFSR) \\
\texttt{/service-load-sr-types}, \texttt{/service-load-heads} & SR Type Master (Service
Load) \\
\texttt{/lead-categories}, \texttt{/lead-cats} & Lead Category Master \\
\texttt{/se-uid} (+ \texttt{/sync}, \texttt{/import}) & SE UID Master (Master Admin only) \\
\texttt{POST /upload}, \texttt{GET /upload/progress}, \texttt{GET /uploads} & the Spare /
Labour import and its progress \\
\texttt{GET /data/summary}, \texttt{/data/preview}, \texttt{POST /data/cancel-row},
\texttt{DELETE /data} & stored data, preview and cancellation \\
\texttt{POST /training/upload}, \texttt{GET /training/report}, \texttt{/training/summary},
\texttt{DELETE /training/data} & Training Report \\
\texttt{GET /report} & Sales and Labour -- All / Branch tables \\
\texttt{GET /report/fy-summary} & FY / Quarterly / Month-wise \\
\texttt{GET /report/branch-detail} & Branch-wise Report \\
\texttt{GET /report/employee-productivity} & Employee Productivity \\
\texttt{GET /report/sr-allocation} & SR Allocation \\
\texttt{GET /report/annual/service-penetration} & Service Penetration \\
\texttt{GET /report/annual/cdi} & Customer Delight Index \\
\texttt{GET /report/annual/amc-bandhan} & AMC and Bandhan Projection \\
\texttt{GET /report/annual/amc-monthly} & the AMC monthly sheet \\
\texttt{GET /report/annual/service-load} & Service Load and Response \\
\texttt{GET/POST /cdi-targets}, \texttt{/amc-targets}, \texttt{/service-load-targets},
\texttt{/service-load-manual} & the AOP figures those sheets are measured against \\
\end{longtable}

\vspace{12pt}
\begin{center}\small\itshape End of document\end{center}

\end{document}
